% OSF pre-registration draft — must be finalized and registered BEFORE
% touching real data in the SCB TRE (see notes/continuity.md).
% Compile with pdflatex (Overleaf-compatible).
\documentclass[11pt,a4paper]{article}

\usepackage[utf8]{inputenc}
\usepackage[T1]{fontenc}
\usepackage[english]{babel}
\usepackage[margin=2.5cm]{geometry}
\usepackage{booktabs}
\usepackage{amsmath}
\usepackage{xcolor}
\usepackage[hidelinks]{hyperref}

\newcommand{\todo}[1]{\textcolor{red}{[TODO: #1]}}

\title{Pre-registration: Digital lifestyles and consumption-based carbon
emissions}
\author{David Andersson \and Mathias Lehner}
\date{Draft --- \today}

\begin{document}

\maketitle

\section{Study information}

\subsection{Research questions}
\begin{enumerate}
  \item \textbf{Net gradient.} What is the association between digital lifestyle
    intensity and total consumption-based CO\textsubscript{2}e, expressed per
    additional daily hour of screen time, controlling for sociodemographic
    factors?
  \item \textbf{Decomposition.} Is the gradient negative in transport
    (demobilization) and positive in goods, e-commerce-intensive categories,
    and digital services (rebound and direct effects), with no-channel
    categories (rent, insurance) as placebos?
  \item \textbf{Mechanism.} Does lower out-of-home leisure frequency among the
    digitally intensive account for part of the transport gradient?
\end{enumerate}

\subsection{Hypotheses}
\todo{State directional hypotheses per research question, incl.\ the placebo
prediction (null gradient in rent and insurance).}

\section{Data description}

This study reuses existing data (the Konsumtionskollen sample): approximately
6{,}000 Swedish adults who completed surveys in 2024 and connected their bank
accounts to a research app that retrieved transactions back to 2019, categorized
them to COICOP, and converted them to CO\textsubscript{2}e. Register linkage
provides income, education, household composition, and DeSO-level urbanity.
Analyses of real data take place in the SCB Trusted Research Environment (TRE).

\paragraph{Prior data access.} \todo{Disclose: analysis code developed against a
synthetic/mock version of the data; the endline survey and real values have not
been analyzed. State precisely what has and has not been seen.}

\section{Variables}

\subsection{Exposure: digital lifestyle intensity index}
A pre-registered index combining:
\begin{itemize}
  \item \textbf{Device-use frequencies} (baseline q11\_1--q11\_5+; 1--6 scale):
    computer, smartphone, TV, game console, tablet.
  \item \textbf{Digital-activity frequencies} (baseline q11b\_1--q11b\_8; 1--6
    scale): streaming, calls, social media, information search, gaming, digital
    reading, e-commerce, online selling.
\end{itemize}
The index is anchored cardinally by the device-assisted screen-time report
(endline E4: average daily screen time in hours:minutes for the previous week,
read from iOS Sk\"armtid / Android Digitalt v\"alm\aa ende).
\todo{Specify the aggregation rule (e.g., standardized mean), the anchoring
regression, reliability reporting ($\alpha$), and handling of E4 opt-outs.}

\subsection{Outcomes}
Person-level annual consumption-based CO\textsubscript{2}e: total, and by
category group (transport; goods; e-commerce-intensive; digital services;
placebo: rent, insurance). \todo{Finalize and justify the COICOP leaf-category
mapping (draft in 00\_constants.R).}

\subsection{Mediator}
Out-of-home leisure frequency (endline E1, 7-point, 12-month recall).

\subsection{Covariates}
Age, gender, income, education, household composition, DeSO-level urbanity
(register-based).

\subsection{Moderators (pre-specified heterogeneity)}
Age, gender, urbanity.

\section{Analysis plan}

\subsection{Models}
\todo{OLS with HC3 robust standard errors; outcome winsorized at P99;
specification per research question, incl.\ the mediation approach for the
mechanism analysis.}

\subsection{Inference criteria}
\todo{Two-sided tests, $\alpha = 0.05$; multiple-comparison handling for the
category decomposition.}

\subsection{Data exclusions}
\todo{E4 opt-outs / missing screen time; transaction-coverage thresholds;
attention check (endline C11)?}

\subsection{Missing data}
\todo{Handling of item non-response in index components and covariates.}

\subsection{Robustness}
\todo{Alternative index weightings; unanchored (ordinal) index; alternative
winsorization; timing sensitivity (screen time measured at endline vs.\
consumption 2019--2024).}

\end{document}
